\documentclass[12pt]{article}

\usepackage[a4paper, margin=1in]{geometry}
\usepackage{amsmath, amssymb, amsthm}
\usepackage{physics}
\usepackage{graphicx}
\usepackage{hyperref}
\usepackage{enumitem}

\title{CSE400 -- Fundamentals of Probability in Computing}
\author{Lecture 11--12: Transformation of Random Variables}
\date{}

\begin{document}

\maketitle

\section{Transformation of Random Variables}

\subsection*{Goal}

Given a random variable $X$ with known PDF $f_X(x)$ and CDF $F_X(x)$, and a transformation

\[
Y = g(X)
\]

find the distribution of $Y$: $F_Y(y)$ and $f_Y(y)$.

\subsection{Monotonic Transformation}

Assume $g(x)$ is invertible.

\subsubsection*{Step 1: CDF Method}

\[
F_Y(y) = P(Y \le y)
\]

Since $Y=g(X)$:

\[
F_Y(y) = P(g(X) \le y)
\]

If $g$ is increasing:

\[
F_Y(y) = P(X \le g^{-1}(y)) = F_X(g^{-1}(y))
\]

If $g$ is decreasing:

\[
F_Y(y) = P(X \ge g^{-1}(y)) = 1 - F_X(g^{-1}(y))
\]

\subsubsection*{Step 2: PDF via Differentiation}

\[
f_Y(y) = \frac{d}{dy}F_Y(y)
\]

For invertible $g$:

\[
f_Y(y) = f_X(g^{-1}(y)) \left| \frac{dx}{dy} \right|
\]

where

\[
x = g^{-1}(y)
\]

\subsubsection*{Notes}

\begin{itemize}
\item Limits must be transformed appropriately.
\item Absolute value handles derivative sign.
\item Applies when $f_X(x)$ is known a priori.
\end{itemize}

\section{Worked Example: Nonlinear Transformation}

\subsection*{Given}

\[
X \sim \text{Uniform}(-1,1)
\]

\[
f_X(x)=
\begin{cases}
\frac12, & -1<x<1 \\
0, & \text{otherwise}
\end{cases}
\]

Transformation:

\[
Y = \sin\left(\frac{\pi X}{2}\right)
\]

\subsection*{Step 1: Inverse Transformation}

\[
x = \frac{2}{\pi}\sin^{-1}(y)
\]

\subsection*{Step 2: Derivative}

\[
\frac{dx}{dy} = \frac{2}{\pi}\frac{1}{\sqrt{1-y^2}}
\]

\subsection*{Step 3: Apply Formula}

\[
f_Y(y) = f_X(x)\left|\frac{dx}{dy}\right|
\]

\[
= \frac12 \cdot \frac{2}{\pi}\frac{1}{\sqrt{1-y^2}}
\]

\[
\boxed{
f_Y(y)=\frac{1}{\pi\sqrt{1-y^2}}, \quad -1<y<1
}
\]

\section{Function of Two Random Variables}

Let

\[
Z = g(X,Y)
\]

Examples:

\begin{itemize}
\item $Z = X + Y$
\item $Z = X - Y$
\item $Z = \dfrac{X}{Y}$
\item $R = \sqrt{X^2 + Y^2}$
\end{itemize}

\section{Distribution of the Sum: $Z=X+Y$}

\subsection*{CDF Approach}

\[
F_Z(z) = P(Z \le z)
\]

\[
= P(X+Y \le z)
\]

\[
= \iint_{x+y \le z} f_{XY}(x,y)\,dx\,dy
\]

\subsection*{PDF via Differentiation}

\[
f_Z(z)=\int_{-\infty}^{\infty} f_{XY}(x,z-x)\,dx
\]

\section{Convolution Formula (Independent Case)}

If $X$ and $Y$ are independent,

\[
f_{XY}(x,y)=f_X(x)f_Y(y)
\]

then

\[
\boxed{
f_Z(z)=\int_{-\infty}^{\infty} f_X(x)\,f_Y(z-x)\,dx
}
\]

This is the \textbf{convolution integral}.

\section{Special Case: Sum of Normal Random Variables}

Given

\[
X \sim N(0,1), \quad Y \sim N(0,1)
\]

independent,

\[
f_Z(z)=\int_{-\infty}^{\infty}
\frac{1}{\sqrt{2\pi}}e^{-x^2/2}
\frac{1}{\sqrt{2\pi}}e^{-(z-x)^2/2}dx
\]

After simplification,

\[
\boxed{Z \sim N(0,2)}
\]

\section{Sum of Exponential Random Variables}

If

\[
f_X(x)=\lambda e^{-\lambda x}, \quad x\ge0
\]
\[
f_Y(y)=\lambda e^{-\lambda y}, \quad y\ge0
\]

then

\[
f_Z(z)=\int_0^z \lambda e^{-\lambda x}\lambda e^{-\lambda(z-x)} dx
\]

\[
= \lambda^2 e^{-\lambda z}\int_0^z dx
\]

\[
\boxed{
f_Z(z)=\lambda^2 z e^{-\lambda z}, \quad z\ge0
}
\]

\[
f_Z(z)=0 \quad \text{otherwise}
\]

\section{Difference of Random Variables}

For

\[
Z = X - Y
\]

\[
F_Z(z)=P(X-Y\le z)
\]

\[
= \iint_{x-y \le z} f_{XY}(x,y)\,dx\,dy
\]

PDF is obtained by differentiation.

\section{Ratio Transformation}

For

\[
Z=\frac{X}{Y}
\]

analysis depends on the sign of $Y$ (cases $Y>0$ and $Y<0$).

\section{Radial Transformation Example}

\[
R=\sqrt{X^2+Y^2}
\]

Find $f_R(r)$.

\section{Leibniz Rule}

If

\[
G(x)=\int_{a(x)}^{b(x)} h(x,y)\,dy
\]

then

\[
\frac{dG(x)}{dx}
= h(x,b(x))\frac{db}{dx}
- h(x,a(x))\frac{da}{dx}
+ \int_{a(x)}^{b(x)} \frac{\partial h(x,y)}{\partial x} dy
\]

\section{Key Takeaways}

\begin{itemize}
\item Transformations derive distributions of new random variables.
\item CDF method is primary; PDF obtained by differentiation.
\item Monotonic transformations use inverse and Jacobian.
\item Functions of two RVs require region integration.
\item Sum of independent RVs $\rightarrow$ convolution.
\item Normal sums remain normal.
\item Exponential sums produce Gamma-type distribution.
\end{itemize}

\end{document}
